\chapter{Symbolic molecular simulation}
\ChapterGoals{%
\item decide what counts as a state of a molecular system, and canonicalise it;
\item implement reactions with preconditions as well as count changes;
\item use stoichiometry as a necessary (not sufficient) reachability test;
\item explore the reachable state space by breadth-first search with an honest stopping contract;
\item state what a found pathway proves, and test a simulator without letting it verify itself.}
{Chapters~20 and~26; graphs and breadth-first search (Appendix~\gref{app:math}).}
A simulator is an executable scientific claim. Before asking how fast a molecular
device computes, ask which transitions its model permits, what it counts as a
state, and what a successful search actually proves. This chapter builds a small,
exact multiset engine and then tries to falsify it with independently calculated
transitions, conservation laws, and path enumeration.

The engine is intentionally not a strand-displacement simulator. It has no domain
matching, geometry, rates, diffusion, energy, or fluorescence. Its purpose is to
separate a precise formal question from the later questions of chemical realization
and experimental observation. Chapter 26's provenance discipline carries forward:
the symbol table, rule set, bounds, and stopping reason belong to the result.

\section{The state is a modeling decision}
Let $\mathcal S=(S_1,\ldots,S_m)$ be an ordered finite species universe. A state is
$x\in\mathbb N^m$, with $x_i$ the exact nonnegative integer count of $S_i$.
Equivalently, it is a finite multiset. The order of dictionary entries has no
physical meaning, so the software stores nonzero counts as an immutable sorted
tuple of species/count pairs.

\Plate{01-ladder}{A conceptual molecular abstraction, not a specified chemical
reaction: strands and complexes are mapped to species identities; hidden topology
is information deliberately discarded by a count model.}{dnad27:ladder}

A symbol may represent a free strand, an entire complex, a tile, or a logical
signal, but these interpretations must not be mixed silently. An $AB$ complex is
one object in the count vector even if its material inventory contains two
constituents. Two molecules with identical sequences but different binding
partners may need different species identities. Tile positions and exposed glues
usually need a spatial configuration, not just a bag of tile counts.

This is an abstraction map $\alpha:\mathcal X_{\rm physical}\to\mathbb N^m$.
If two physical configurations map to the same $x$ but have different legal
transitions, the model has forgotten a relevant precondition. One remedy is to
refine the species universe; another is to use a graph-based complex state. Neither
is equivalent to adding a guessed rate constant to an inadequate state.

\section{Canonicalization is part of correctness}
\Listing{canon}
The implementation rejects fractional, negative, and Boolean counts. A Boolean is
not accepted as an integer count even though Python treats \texttt{bool} as an
integer subtype. Duplicate species identities in a tuple are rejected rather than
merged silently. Zero counts disappear from storage, and unknown species are
rejected when a state enters a declared network.

These decisions prevent three different failure modes: truncating $1.9$ molecules
to one, treating two encodings of the same multiset as distinct states, and creating
an undeclared species by a spelling error. Canonicalization is not a chemistry
validator. It only establishes an unambiguous finite-multiset representation.

\section{A reaction has a precondition, not just a delta}
A rule $r$ contains a unique label, consumed vector $C_r$, and produced vector
$P_r$. Its operational semantics is
\[
x\xrightarrow{r}x'\quad\Longleftrightarrow\quad
x\ge C_r,\qquad x'=x-C_r+P_r.
\]
The inequality is componentwise. Consumption and production are one atomic rewrite,
not a sequence of partly visible removals and additions.

\Plate{02-rule}{Consumed and produced inventories are distinct from their net
difference. A catalyst can have zero net change and still be an essential
enabledness precondition.}{dnad27:rule}
\Listing{fire}

For $2A+E\to B+E$, a state with three $A$ and one $E$ becomes $A+B+E$.
Without $E$ the rule is disabled, even though the net increment is merely
$-2A+B$. Testing only whether the resulting vector is nonnegative would wrongly
permit an enzyme-free transition. This catalyst example is a decisive
counterexample to delta-only rule implementations.

The reference accepts empty consumed or produced multisets. Such rules can model
abstract input or removal, but an empty-consumption production rule can create an
unbounded state space. The software allows the formal construct; the modeler must
justify its interpretation and cannot call it closed material conservation.

\section{Stoichiometry gives necessary constraints}
Arrange rule increments as columns of the stoichiometric matrix
\[
S=[P_1-C_1\ \cdots\ P_q-C_q]\in\mathbb Z^{m\times q}.
\]
After a legal path whose firing-count vector is $\sigma\in\mathbb N^q$,
\[
x=x_0+S\sigma.
\]
This \Term{state equation} is necessary, not sufficient. Consider $B\to B+D$
starting from one $A$ and no $B$. The candidate $A+D=x_0+S$ is nonnegative, yet no
rule is enabled because the unchanged $B$ catalyst is absent. A linear equation
cannot recover the order-sensitive applicability conditions it discarded.

\Listing{stoichiometry}
For the assembly rules $A+B\to AB$ and $AB+C\to ABC$, with species order
$(A,B,C,AB,ABC)$, the matrix is
\[
S=\begin{pmatrix}
-1&0\\-1&0\\0&-1\\1&-1\\0&1
\end{pmatrix}.
\]
For $w=(1,1,1,2,3)^\top$, direct multiplication gives $w^\top S=(0,0)$.
Consequently $w^\top x=3$ along every legal path from $A+B+C$.

\Plate{05-invariant}{An independently derived left-nullspace invariant checks
material-unit accounting. It cannot establish enabledness, timing, or experimental
mass conservation by itself.}{dnad27:invariant}

\begin{proposition}
If $w^\top(P_r-C_r)=0$ for every rule, then $w^\top x$ is invariant on every
finite legal path.
\end{proposition}
\begin{proof}
For one transition, $w^\top x'=w^\top x+w^\top(P_r-C_r)=w^\top x$.
Induction on path length proves the claim.
\end{proof}
With strictly positive integer weights and fixed total $W$, each count is at most
$\lfloor W/w_i\rfloor$. Thus conservation can establish a finite reachable state
space, bounded above by $\prod_i(1+\lfloor W/w_i\rfloor)$ configurations. This
bound is often loose; it is a finiteness argument, not a predicted graph size.

\section{States merge; labeled transitions do not}
A labeled transition system has states as vertices and triples $(x,r,x')$ as
edges. Different rules can produce the same successor, and a rule can return to
the same state. Both distinctions matter when diagnosing competing mechanisms.

\Plate{03-graph}{Two assembly routes share a final state. State deduplication merges
that vertex but retains all four labeled reaction edges. A first-discovery witness
records one shortest route, not all routes.}{dnad27:graph}

Add $A+C\to AC$ and $AC+B\to ABC$. Starting from $A+B+C$, the resulting diamond
has four states and four labeled edges. Recording only a predecessor tree would
keep three edges and lose the alternative route. Our search therefore stores a
first-discovery predecessor for witness reconstruction and a separate labeled
edge collection for expanded states.

A second fixture has $A\to B$ and $A\to C$. Its two terminal states are a
counterexample to unique terminal outcome. A third fixture adds an $A$ self-loop
and a $B\to A$ return edge. The finite graph closes, but an execution can loop
forever. Graph closure does not imply termination.

\section{Breadth-first search with an honest stopping contract}
\Listing{bounded_reachable}
Breadth-first search assigns the shortest discovered distance $d(x)$ from $x_0$.
For every state with $d(x)<K$, it enumerates every enabled rule. Canonical equality
merges only identical count states. On first discovery it records the parent and
rule label, then queues the state.

\Plate{04-bfs}{Three stopping outcomes have different logical meanings:
depth-complete exploration, a fully closed reachable set, and an interrupted
state-budget search. An unknown boundary count is not reported as zero.}{dnad27:bfs}

The output separates \texttt{complete\_within\_depth} from
\texttt{exhaustive\_closure}. At depth $K$, successor targets are inspected without
being added to the stored graph. If any target is outside the discovered set,
the result is depth-bounded. If none is outside and all queued states were
processed, the discovered set is closed under every rule and contains the whole
reachable set. This includes cyclic closed graphs.

The reported edge count contains edges expanded from states with $d<K$.
Boundary edges are inspected for closure but are not inserted into that count.
This convention must accompany any table: at $K=0$, a self-loop can establish
closure while the reported expanded-edge count remains zero.

\begin{center}\input{\Generated/trace.tex}\end{center}
Here the initial state is $A+B+C$ and only the two baseline assembly rules are
present. At $K=1$, $ABC$ is absent but reachable just beyond the bound. At $K=2$,
the three-state graph closes and the single terminal state is $ABC$. A larger
bound changes nothing.

A state budget is a separate resource constraint. When the next new state would
exceed it, the search returns \texttt{state\_budget\_exhausted},
both completeness flags false, and an unknown outside-successor count represented
by \texttt{None}. Some discovered states have not yet been expanded; even the
terminal list may be partial. Existing positive witnesses remain useful, but
negative answers are undecided.

\section{What a witness proves}
\Listing{witness}
Following first-discovery predecessors backward and reversing their rule labels
gives a shortest path. For the baseline target, the result is
\texttt{bind\_AB, bind\_C}. The zero-step initial-state witness is empty.

A witness is a constructive reachability certificate only after its labels are
replayed against the declared rules and initial state. The helper assumes the
search's internal predecessor map; it rejects cycles but does not authenticate
an externally supplied map. Archive the symbol table, rules, initial state,
target, search status, and replay trace together.

\begin{proposition}
A completed, unpruned breadth-first search is sound for the declared multiset
semantics and complete for targets at shortest distance at most $K$.
\end{proposition}
\begin{proof}
Each generated edge passes the exact consumed-count precondition and applies the
declared increment, proving soundness. Induct on distance: the initial state is
stored; every successor of a stored state at distance below $K$ is either newly
queued or already represented by the identical canonical state. Therefore every
state at distance at most $K$ is stored. A resource interruption invalidates the
inductive completeness claim.
\end{proof}

\section{Testing without asking the implementation to verify itself}
The independent oracle enumerates path frontiers using dense integer vectors,
not the canonicalizer, firing function, or deduplicating BFS. For all 27 initial
triples with $A,B,C$ counts in $\{0,1,2\}$ and bounds zero through three, it compares
the full state-to-shortest-depth mapping. Every stored edge is also replayed with
independent consumed/produced vector arithmetic.

Further tests retain parallel labels and self-loops, reject undeclared species,
validate multiplicity and catalysts, detect missing invariant weights, exercise
budget interruption, and distinguish multiple terminals from acceptance. The
left-nullspace test multiplies the hand-written matrix by the hand-derived
weights; it does not reuse the invariant helper as its oracle.

These tests establish properties of this implementation and its finite fixtures.
They are not a proof for arbitrary richer molecular languages. Partial-order
reduction, symmetry quotienting, or approximate merging requires a separate
argument that the property being queried is preserved.

\section{Acceptance, confluence, and fairness are separate}
An acceptance predicate $A(x)$ is supplied by the application. Terminality means
no rule is enabled; neither condition implies the other. An accepted state can
still react, and an unwanted dead end can be terminal.

For a closed finite graph, one can inspect whether every reachable terminal is
accepted and whether all reachable states can reach an accepted state. Neither
property guarantees that every execution eventually accepts: a scheduler may
remain on a cycle. Universal termination additionally requires excluding reachable
cycles, or specifying and proving a suitable fairness condition.

Confluence is stronger than observing one successful path. Every pair of reachable
branches must admit a common descendant under the declared equivalence. A shared
terminal diamond is an illustrative join, not a proof that a larger rule set is
globally confluent. Strongly connected components and graph reachability algorithms
can extend the reference, but must keep their logical questions explicit.

\section{Formal possibility is not chemical realizability}
\Plate{06-boundary}{The evidence boundary: formal transitions establish possible
paths; kinetic models add distributions and time; experiments test whether the
specified molecules and readout support those models.}{dnad27:boundary}
Cardelli's author-hosted discussion of molecular semantics distinguishes models
of interactions from quantitative rate interpretations. The reference here is
a count rewrite system, not an implementation of an entire process calculus
\cite{Cardelli2008}.

DNA strand-displacement constructions can approximate prescribed reaction-network
dynamics under specified designs and scaling conditions \cite{Soloveichik2010}.
Auxiliary species, rate rescaling, leakage, and implementation limitations do not
vanish under that correspondence. Our $AB$ drawing is a conceptual assembly
inventory, not a validated domain-level construction.

\begin{KeyIdeas}
\begin{itemize}\setlength{\itemsep}{2pt}
\item Choosing the state (which species, counted how) is a modelling decision;
canonicalisation makes equal states compare equal.
\item A reaction has a precondition (enough reactants) as well as a change
vector; the change vector alone permits impossible moves.
\item Stoichiometric feasibility is necessary for reachability but not
sufficient.
\item Breadth-first search over a bounded state space finds shortest pathways;
a cutoff must be reported as a cutoff, not as closure.
\item A witness pathway proves that an outcome is possible in the model;
acceptance, confluence and fairness are further properties.
\item Formal possibility is not chemical realisability; rates, energies and
geometry come next (Chapter~28).
\end{itemize}
\end{KeyIdeas}

\section{Worked problems and design exercises}
\begin{enumerate}
\item Canonicalize $2A+0B+C$. \textbf{Answer:} $((A,2),(C,1))$.
\item Why reject count $1.9$? \textbf{Answer:} it is not a discrete molecule count;
truncation changes the model.
\item Enable $2A+E\to B+E$ in $A+E$? \textbf{Answer:} no; one $A$ is missing.
\item Fire it from $3A+E$. \textbf{Answer:} $A+B+E$.
\item Why does $E$'s zero matrix row not remove its precondition?
\textbf{Answer:} consumed and produced requirements are not encoded by their difference.
\item Derive the assembly matrix. \textbf{Answer:} subtract each consumed vector
from its produced vector in the stated species order.
\item Verify the left-nullspace condition for $w=(1,1,1,2,3)^\top$. \textbf{Answer:}
$w^\top S=(0,0)$.
\item Bound $x_{ABC}$ at total weight three. \textbf{Answer:} at most one.
\item Does matching total weight establish reachability? \textbf{Answer:} no;
it is only a necessary invariant constraint.
\item Give a feasible state-equation target that is unreachable.
\textbf{Answer:} $A+D$ under $B\to B+D$ from $A$.
\item Count the diamond states and edges. \textbf{Answer:} four of each.
\item What does a predecessor tree lose? \textbf{Answer:} alternative labels,
non-tree edges, and cycles.
\item Interpret no $ABC$ at bound one. \textbf{Answer:} absent within one step,
not globally unreachable.
\item When is a bounded result nevertheless exhaustive? \textbf{Answer:} all
queued states processed and every boundary successor already inside the set.
\item What does a budget interruption do to negative answers?
\textbf{Answer:} it leaves them undecided.
\item Why can closure coexist with nontermination? \textbf{Answer:} a reachable
cycle admits an infinite execution.
\item Is an accepted state necessarily terminal? \textbf{Answer:} no; acceptance
is a separate predicate.
\item Design a richer strand state. \textbf{Rubric:} domains, strand orientation,
complex membership, binding topology, and explicit applicability rules.
\item Propose symmetry reduction. \textbf{Rubric:} define an equivalence and prove
it preserves the acceptance/reachability question before merging states.
\item Design a replay archive. \textbf{Rubric:} versioned universe/rules, initial
state, bounds/budget, complete status, target, labeled witness, independent replay,
and the molecular abstraction's omissions.
\end{enumerate}

\section{Bridge to quantitative simulation}
Chapter 28 adds time and rates through deterministic and stochastic semantics.
Ignoring timing should recover these same consumed/produced transitions. Agreement
on stoichiometry is the first gate, not the last: units, propensities, stiffness,
finite-copy noise, and readout adequacy require new independent tests.
